\section{Discussion}
\label{sec:discussion}

% TODO: Interpret results in light of the three RQs once data are available.

\subsection{Context Breadth as a Design Variable (RQ1)}
\label{sec:discussion:breadth}

% TODO: Discuss the C1 → C3 gradient and what results imply about the marginal
%   value of each additional surface in the agent's context. Connect to the
%   information-foraging literature. Note the scope boundary: this speaks to how
%   much context an embedded agent needs, NOT to all-in-one vs. separate tools.

\subsection{Task Type as a Moderator (H2)}
\label{sec:discussion:tasktype}

% TODO: Discuss whether cross-surface tasks benefit disproportionately from
%   broader context (as hypothesized). When is a narrower context sufficient?
%   T6 (GitHub-only-identifiable PR) is the cleanest top-of-ladder test.

\subsection{Reliance and Division of Labor (RQ2)}
\label{sec:discussion:reliance}

% TODO: Discuss how delegation shifts with context breadth and whether reliance
%   is calibrated — do users verify/override where the agent is actually wrong,
%   or under/over-rely? Connect to validation-and-repair behavior in LLM-assisted
%   work.

\subsection{Trust versus Reality (RQ3 / H5)}
\label{sec:discussion:trust}

% TODO: Discuss the perception–behavior gap — whether trust and perceived
%   speedup track measured effectiveness or diverge (e.g., agents that *do more*
%   trusted even when no faster). This is the most novel angle; flag design
%   implications of mis-calibrated trust in action-capable agents.

\subsection{Code-Enforced Context Scoping as a Design Principle}
\label{sec:discussion:gate}

The \ContextGate pattern---enforcing context boundaries in server code rather
than model prompts---has implications beyond this study. First, it makes the
manipulation mechanically verifiable: any reviewer can inspect the
\texttt{agent\_context\_sent} audit logs and confirm that a \CI\ agent never
received board data. Second, it suggests a governance model for production AI
agents in collaborative tools, where organizations need auditable controls over
what information agents can observe. Third, it decouples the \emph{scope} of an
agent from its \emph{capability}, allowing researchers and practitioners to vary
context breadth independently of model capability.

\subsection{Limitations}
\label{sec:discussion:limitations}

% TODO: Address:
%   - Controlled lab setting vs. naturalistic use
%   - Seeded PetTracker content may not generalize to all domains
%   - Single session; no longitudinal effects measured
%   - Between-subjects design: no within-person learning-curve data across conditions
%   - Statistical power (dependent on final N)
%   - Generalization to other LLM models (study used a single pinned model)
